\documentclass[11pt,oneside]{article}
\usepackage[a4paper,margin=27mm,headheight=15pt]{geometry}
\usepackage[T1]{fontenc}
\usepackage[utf8]{inputenc}
\IfFileExists{libertinus.sty}{\usepackage{libertinus}}{\usepackage{lmodern}}
\usepackage{microtype}
\usepackage{amsmath,amssymb,amsthm,mathtools,mathrsfs}
\usepackage{bm}
\usepackage{booktabs,longtable,array,tabularx}
\usepackage{graphicx}
\usepackage{pgfplots}
\usepackage{xcolor}
\usepackage{enumitem}
\usepackage{fancyhdr}
\usepackage{titlesec}
\usepackage{listings}
\usepackage{xurl}
\usepackage{hyperref}
\usepackage[nameinlink,noabbrev]{cleveref}

\definecolor{linkblue}{RGB}{25,75,135}
\definecolor{shadegray}{RGB}{246,247,249}
\definecolor{rulegray}{RGB}{195,201,208}
\pgfplotsset{compat=1.18}
\hypersetup{
  colorlinks=true,
  linkcolor=linkblue,
  citecolor=linkblue,
  urlcolor=linkblue,
  pdftitle={A Formula Atlas for the Thue-Morse Sequence},
  pdfsubject={Binary digits, valuations, Pascal support, products, transforms, and Prouhet cancellation},
  pdfkeywords={Thue-Morse, binary digit sum, Pascal triangle, Prouhet, Mahler equation, Walsh transform, Riesz product}
}

\pagestyle{fancy}
\fancyhf{}
\fancyhead[L]{\small Thue--Morse formula atlas}
\fancyhead[R]{\small\nouppercase{\leftmark}}
\fancyfoot[C]{\thepage}
\renewcommand{\headrulewidth}{0.4pt}

\titleformat{\section}{\Large\bfseries}{\thesection}{0.7em}{}
\titleformat{\subsection}{\large\bfseries}{\thesubsection}{0.7em}{}
\titleformat{\subsubsection}{\normalsize\bfseries}{\thesubsubsection}{0.7em}{}
\setcounter{secnumdepth}{3}
\setcounter{tocdepth}{2}
\setlist[itemize]{leftmargin=2em,itemsep=0.25em,topsep=0.4em}
\setlist[enumerate]{leftmargin=2.3em,itemsep=0.25em,topsep=0.4em}
\setlength{\emergencystretch}{2em}

\newtheorem{theorem}{Theorem}[section]
\newtheorem{proposition}[theorem]{Proposition}
\newtheorem{lemma}[theorem]{Lemma}
\newtheorem{corollary}[theorem]{Corollary}
\newtheorem{conjecture}[theorem]{Conjecture}
\theoremstyle{definition}
\newtheorem{definition}[theorem]{Definition}
\newtheorem{algorithm}[theorem]{Algorithm}
\newtheorem{example}[theorem]{Example}
\theoremstyle{remark}
\newtheorem{remark}[theorem]{Remark}
\newtheorem{warning}[theorem]{Warning}

\newcommand{\R}{\mathbb R}
\newcommand{\C}{\mathbb C}
\newcommand{\Q}{\mathbb Q}
\newcommand{\N}{\mathbb N}
\newcommand{\Z}{\mathbb Z}
\newcommand{\F}{\mathbb F}
\newcommand{\wt}{\operatorname{wt}_2}
\newcommand{\swt}{s_2}
\newcommand{\vtwo}{v_2}
\newcommand{\e}{\mathrm e}
\newcommand{\ii}{\mathrm i}
\newcommand{\dd}{\,\mathrm d}
\newcommand{\sgn}{\operatorname{sgn}}
\newcommand{\eps}{\varepsilon}
\newcommand{\supp}{\operatorname{supp}}
\newcommand{\xor}{\mathbin{\mathsf{xor}}}
\newcommand{\submask}{\preccurlyeq}
\newcommand{\repo}{\nolinkurl{Analysis/FabiusFunction}}
\newcommand{\lean}[1]{%
  \ifmmode\text{\nolinkurl{#1}}\else\nolinkurl{#1}\fi}
\newcommand{\fall}[2]{#1^{\underline{#2}}}
\newcommand{\Cat}{\operatorname{Cat}}
\newcommand{\ord}{\operatorname{ord}}
\newcommand{\1}{\mathbf 1}

\newenvironment{proofidea}{\par\smallskip\noindent\textit{Proof architecture.}\ }{\hfill$\square$\par\smallskip}
\newenvironment{boxedremark}{%
  \par\smallskip\noindent\begin{center}\begin{minipage}{0.94\textwidth}%
  \color{black}\hrule\vspace{0.6em}\small
}{%
  \vspace{0.6em}\hrule\end{minipage}\end{center}\smallskip
}

\lstdefinestyle{pseudo}{
  basicstyle=\ttfamily\small,
  backgroundcolor=\color{shadegray},
  frame=single,
  rulecolor=\color{rulegray},
  columns=fullflexible,
  keepspaces=true,
  showstringspaces=false,
  xleftmargin=0.7em,
  xrightmargin=0.7em,
  aboveskip=0.8em,
  belowskip=0.8em
}

\title{\bfseries A Formula Atlas for the Thue--Morse Sequence\\[0.35em]
\large Binary digits, valuations, Pascal support, products, transforms, and Prouhet cancellation}
\author{Companion report to the Fabius--Rvachev formal development\\
\small Vladimir Reshetnikov's \texttt{ProveIt} repository}
\date{Repository snapshot: 25 August 2026\\
\small based on commit \nolinkurl{eaea1b9afe2b7ff0b7dd880bd1710e303dec6d80}}

\begin{document}
\maketitle

\begin{abstract}
Section 11.8 of the repository exposition derives a striking formula for the zero--one
Thue--Morse bit by counting odd entries in one row of Pascal's triangle and then taking an
exact base-two logarithm.  The purpose of this report is to build a much wider formula atlas
around that observation.  We pass systematically among the two normalizations
\[
  \tau(n)=\swt(n)\bmod 2,\qquad
  \eps_n=(-1)^{\swt(n)}=1-2\tau(n),
\]
and derive self-contained formulas from binary digits, floor and fractional-part sums,
trigonometric signs, factorial valuations, central binomial coefficients, Catalan numbers,
carry counts, multinomial parity, finite and infinite generating products, Mahler equations,
ordinary Fourier inversion, Walsh characters, finite-difference operators, Prouhet
cancellation, exact higher moments, and formal power series over \(\F_2\).  Two particularly
useful extensions of the Pascal-row viewpoint are developed in detail: an \(r\)-nomial
logarithm family in which the number of odd multinomial coefficients is
\(r^{\swt(n)}\), and a sparse Prouhet theorem on the odd support of an arbitrary Pascal row.
We also give exact prefix-discrepancy formulas, autocorrelation recurrences, the finite Riesz
products of the Thue--Morse spectral measure, a classical evaluated infinite product, and a
base-\(q\) perspective.

The term ``new'' below always means new relative to the cited subsection or newly derived in
this report; no claim of historical priority is intended.  Except where a literature theorem
is explicitly labelled as such, every displayed identity is proved here.  The existing
binomial--logarithm theorem has named Lean counterparts in the repository; the additional
families in this report are proposed mathematical extensions and are not claimed to be
formalized there.
\end{abstract}

\tableofcontents
\newpage

\section{Scope, conventions, and a first atlas}
\label{sec:scope}

\subsection{The two normalizations}

Throughout, \(\N=\{0,1,2,\ldots\}\).  Write
\[
 n=\sum_{j\ge 0} b_j(n)2^j,
 \qquad b_j(n)\in\{0,1\},
 \qquad \swt(n)=\sum_{j\ge0}b_j(n).
\]
Only finitely many digits are nonzero.  The zero--one and signed Thue--Morse sequences are
\begin{equation}
 \tau(n)=\swt(n)\bmod2\in\{0,1\},
 \qquad
 \eps_n=(-1)^{\swt(n)}=1-2\tau(n)\in\{1,-1\}.
 \label{eq:normalizations}
\end{equation}
Thus every formula for one normalization immediately yields a formula for the other:
\begin{equation}
 \tau(n)=\frac{1-\eps_n}{2},
 \qquad
 \eps_n=1-2\tau(n).
 \label{eq:conversion}
\end{equation}
The first values are
\[
 (\tau(n))_{n\ge0}=0,1,1,0,1,0,0,1,1,0,0,1,0,1,1,0,\ldots,
\]
\[
 (\eps_n)_{n\ge0}=1,-1,-1,1,-1,1,1,-1,-1,1,1,-1,1,-1,-1,1,\ldots.
\]

\begin{boxedremark}
The signed normalization is multiplicative under binary concatenation and factors cleanly in
generating functions.  The zero--one normalization is the natural parity bit and is the one
used in the binomial--logarithm theorem.  Keeping both visible prevents avoidable sign errors.
\end{boxedremark}

\subsection{What section 11.8 already proves}

Let
\begin{equation}
 \mathcal X(n)=\sum_{k=0}^{n}(-1)^{\binom nk}.
 \label{eq:source-X}
\end{equation}
The cited section proves that the number \(O(n)\) of odd entries in row \(n\) of Pascal's
triangle is \(2^{\swt(n)}\), and consequently
\begin{equation}
 n+1-\mathcal X(n)=2O(n)=2^{\swt(n)+1}.
 \label{eq:source-exact-power}
\end{equation}
The logarithm is therefore exact, not numerical, and
\begin{equation}
 \tau(n)
 =\frac{1+(-1)^{\log_2\!\left(n+1-\sum_{k=0}^{n}(-1)^{\binom nk}\right)}}{2}.
 \label{eq:source-formula}
\end{equation}
Our first simplification is to remove the harmless shift by one:
\begin{equation}
 \boxed{
 \eps_n=(-1)^{\log_2 O(n)},
 \qquad
 \tau(n)=\frac{1-(-1)^{\log_2 O(n)}}2.}
 \label{eq:pascal-count-direct}
\end{equation}
The substance remains the same: binary parity is encoded in the support size of a polynomial
modulo two.

\subsection{A one-page formula index}

The following table previews the principal exact representations proved below.  The symbol
\(k\submask n\) means that every one-bit of \(k\) is also a one-bit of \(n\), and
\(P_m(z)=\sum_{0\le n<2^m}\eps_nz^n\).

\renewcommand{\arraystretch}{1.22}
\begin{longtable}{@{}p{0.24\textwidth}p{0.69\textwidth}@{}}
\toprule
Viewpoint & Representative formula \\
\midrule
\endfirsthead
\toprule
Viewpoint & Representative formula \\
\midrule
\endhead
Binary recursion & \(\eps_{2n}=\eps_n\), \(\eps_{2n+1}=-\eps_n\). \\
Binary digits & \(\eps_n=\prod_{j\ge0}(1-2b_j(n))\). \\
Floors & \(\eps_n=(-1)^{\,n-\sum_{j\ge1}\lfloor n/2^j\rfloor}\). \\
Fractional parts & \(\swt(n)=\sum_{j\ge1}\{n/2^j\}\). \\
Trigonometric signs & \(\eps_n=\prod_{j\ge0}\sgn\sin\!\bigl(\pi(n+\tfrac12)/2^j\bigr)\). \\
Factorials & \(\eps_n=(-1)^{n-\vtwo(n!)}\). \\
Central binomial & \(\eps_n=(-1)^{\vtwo\binom{2n}{n}}\). \\
Carries & \(\eps_{m+n}=\eps_m\eps_n(-1)^{\vtwo\binom{m+n}{m}}\). \\
Pascal support & \(O(n)=2^{\swt(n)}\), hence \(\eps_n=(-1)^{\log_2O(n)}\). \\
Multinomial support & \(O_r(n)=r^{\swt(n)}\), hence \(\eps_n=(-1)^{\log_rO_r(n)}\). \\
Finite product & \(P_m(z)=\prod_{j=0}^{m-1}(1-z^{2^j})\). \\
Infinite product & \(\sum_{n\ge0}\eps_nz^n=\prod_{j\ge0}(1-z^{2^j})\), \(|z|<1\). \\
Mahler equation & \(E(z)=(1-z)E(z^2)\). \\
Walsh character & \(\eps_n=(-1)^{\mathbf1\cdot(b_0,\ldots,b_{m-1})}\), \(n<2^m\). \\
Prouhet cancellation & \(\sum_{n<2^m}\eps_np(n)=0\) for \(\deg p<m\). \\
First surviving moment & \(\sum_{n<2^m}\eps_nn^m=(-1)^m m!2^{\binom m2}\). \\
Sparse Pascal support & \(\sum_{k\submask n}\eps_kz^k=\prod_{j:b_j(n)=1}(1-z^{2^j})\). \\
Prefix discrepancy & \(\sum_{n<N}\eps_n=0\) for even \(N\), and \(\eps_{(N-1)/2}\) for odd \(N\). \\
Characteristic two & \((1+x)^3\Theta(x)^2+(1+x)^2\Theta(x)+x=0\). \\
\bottomrule
\end{longtable}
\renewcommand{\arraystretch}{1}

\section{Binary recursion, concatenation, and finite-state formulas}
\label{sec:binary}

\subsection{The fundamental two-scale recurrence}

\begin{proposition}[Even--odd recurrence]
\label{prop:evenodd}
For every \(n\ge0\),
\begin{equation}
 \eps_{2n}=\eps_n,
 \qquad
 \eps_{2n+1}=-\eps_n,
 \label{eq:eps-evenodd}
\end{equation}
and
\begin{equation}
 \tau(2n)=\tau(n),
 \qquad
 \tau(2n+1)=1-\tau(n).
 \label{eq:tau-evenodd}
\end{equation}
Equivalently,
\begin{equation}
 \eps_n=(-1)^n\eps_{\lfloor n/2\rfloor}.
 \label{eq:one-line-recursion}
\end{equation}
\end{proposition}

\begin{proof}
Multiplication by two shifts the binary expansion left and does not change its number of
one-bits.  Adding one to an even number appends a one-bit, increasing the weight by one.
Equations \eqref{eq:eps-evenodd} and \eqref{eq:tau-evenodd} follow.  The final formula is the
two cases combined.
\end{proof}

\begin{corollary}[Complement-and-concatenate construction]
\label{cor:morphism}
Let \(W_m=\tau(0)\tau(1)\cdots\tau(2^m-1)\) be the length-\(2^m\) prefix as a binary word,
and let \(\overline{W_m}\) denote bitwise complementation.  Then
\begin{equation}
 W_{m+1}=W_m\,\overline{W_m}.
 \label{eq:morphism-concat}
\end{equation}
In the signed normalization the second half is the negative of the first.
\end{corollary}

\begin{proof}
For \(0\le r<2^m\), the number \(2^m+r\) has one more leading one-bit than \(r\), so
\(\tau(2^m+r)=1-\tau(r)\) and \(\eps_{2^m+r}=-\eps_r\).
\end{proof}

\subsection{Binary concatenation and the two-element kernel}

\begin{theorem}[No-carry concatenation law]
\label{thm:concat}
For integers \(a,m\ge0\) and \(0\le r<2^m\),
\begin{equation}
 \swt(2^ma+r)=\swt(a)+\swt(r),
 \qquad
 \eps_{2^ma+r}=\eps_a\eps_r,
 \label{eq:concat-sign}
\end{equation}
while
\begin{equation}
 \tau(2^ma+r)=\tau(a)\xor\tau(r).
 \label{eq:concat-bit}
\end{equation}
\end{theorem}

\begin{proof}
The lower \(m\) bits of \(2^ma+r\) are exactly the bits of \(r\), and the remaining bits are
those of \(a\).  There is no overlap and hence no carry.  Additivity of the weight gives the
first formula; exponentiating \((-1)\) and reducing modulo two give the other two.
\end{proof}

\begin{corollary}[The two-kernel]
\label{cor:2kernel}
Every subsequence of the form \((\eps_{2^mn+r})_{n\ge0}\), with \(0\le r<2^m\), is either
\((\eps_n)_{n\ge0}\) or its negative.  Explicitly,
\begin{equation}
 \eps_{2^mn+r}=\eps_r\eps_n.
 \label{eq:kernel}
\end{equation}
Thus the signed sequence has a two-element \(2\)-kernel.
\end{corollary}

\begin{remark}
The recurrence may be read as a two-state automaton: start in state \(+1\), read the binary
digits of \(n\), and flip the state whenever a one is read.  The output state is \(\eps_n\).
This is the most economical finite-state form of the sequence.
\end{remark}

\subsection{Tensor and bitwise-XOR forms}

For \(n<2^m\), let \(\bm b(n)=(b_0(n),\ldots,b_{m-1}(n))\in\F_2^m\) and let
\(\bm1=(1,\ldots,1)\).  Then
\begin{equation}
 \eps_n=(-1)^{\bm1\cdot\bm b(n)}.
 \label{eq:parity-character}
\end{equation}
Consequently the length-\(2^m\) signed block is the Kronecker power
\begin{equation}
 (\eps_0,\ldots,\eps_{2^m-1})=(1,-1)^{\otimes m}.
 \label{eq:tensor}
\end{equation}
This identity is simultaneously a compact formula, a divide-and-conquer construction, and the
reason that the sequence appears as a row of a Walsh--Hadamard matrix.

\begin{algorithm}[XOR-fold parity]
\label{alg:xor-fold}
For a machine word wide enough to hold \(n\), repeatedly XOR the word with a right shift by
powers of two.  The least significant bit of the result is \(\tau(n)\):
\begin{lstlisting}[style=pseudo]
x := n
x := x xor (x >> 1)
x := x xor (x >> 2)
x := x xor (x >> 4)
x := x xor (x >> 8)
... continue while the shift is below the word width ...
tau := x & 1
sign := 1 - 2*tau
\end{lstlisting}
\end{algorithm}

\begin{proofidea}
After the first line, bit zero contains the XOR of the first two input bits; after the shift by
\(2\), it contains the XOR of the first four; after the shift by \(4\), of the first eight;
and so on.  At termination it contains the XOR of all bits, which is their sum modulo two.
\end{proofidea}

\section{Digit, floor, fractional-part, and trigonometric formulas}
\label{sec:digits}

\subsection{Extracting each binary digit arithmetically}

\begin{lemma}[Floor extraction of a bit]
\label{lem:bit-floor}
For every \(j,n\ge0\),
\begin{equation}
 b_j(n)=\left\lfloor\frac{n}{2^j}\right\rfloor
       -2\left\lfloor\frac{n}{2^{j+1}}\right\rfloor.
 \label{eq:bit-floor}
\end{equation}
\end{lemma}

\begin{proof}
Write \(\lfloor n/2^j\rfloor=2q+r\) with \(r\in\{0,1\}\).  Then
\(q=\lfloor n/2^{j+1}\rfloor\), and the remainder \(r\) is precisely the \(j\)-th bit.
\end{proof}

\begin{theorem}[Digit-product formula]
\label{thm:digit-product}
For every \(n\ge0\),
\begin{equation}
 \boxed{\eps_n=\prod_{j\ge0}\bigl(1-2b_j(n)\bigr)}.
 \label{eq:digit-product}
\end{equation}
The product is finite in effect, because all sufficiently high digits vanish.  Therefore
\begin{equation}
 \tau(n)=\frac12\left(1-\prod_{j\ge0}\bigl(1-2b_j(n)\bigr)\right).
 \label{eq:tau-digit-product}
\end{equation}
Substituting \eqref{eq:bit-floor} gives a formula involving only integer arithmetic and floors.
\end{theorem}

\begin{proof}
Each zero-bit contributes \(+1\), each one-bit contributes \(-1\), and their product is
\((-1)^{\sum_jb_j(n)}\).
\end{proof}

Expanding the product gives an exact multilinear polynomial in the bits:
\begin{corollary}[Boolean inclusion--exclusion polynomial]
\label{cor:boolean-poly}
If \(J\) is any finite set containing the support of the binary expansion of \(n\), then
\begin{equation}
 \tau(n)=
 \sum_{\varnothing\ne S\subseteq J}
 (-2)^{|S|-1}\prod_{j\in S}b_j(n).
 \label{eq:boolean-poly}
\end{equation}
\end{corollary}

\begin{proof}
Expand \(\prod_{j\in J}(1-2b_j)\) in \eqref{eq:tau-digit-product} and collect the nonempty
subsets.
\end{proof}

\subsection{Legendre's floor sum and a fractional-part identity}

\begin{proposition}[Floor-sum formula for binary weight]
\label{prop:floor-weight}
For every \(n\ge0\),
\begin{equation}
 \swt(n)=n-\sum_{j\ge1}\left\lfloor\frac{n}{2^j}\right\rfloor.
 \label{eq:weight-floor}
\end{equation}
Consequently
\begin{equation}
 \boxed{
 \eps_n=(-1)^{\,n-\sum_{j\ge1}\lfloor n/2^j\rfloor},
 \qquad
 \tau(n)\equiv n-\sum_{j\ge1}\left\lfloor\frac{n}{2^j}\right\rfloor\pmod2.}
 \label{eq:floor-sign}
\end{equation}
\end{proposition}

\begin{proof}
Starting from \eqref{eq:bit-floor} and summing over \(j\ge0\), the series telescopes:
\[
 \sum_{j\ge0}b_j(n)
 =\sum_{j\ge0}\left(\left\lfloor\frac n{2^j}\right\rfloor
       -2\left\lfloor\frac n{2^{j+1}}\right\rfloor\right)
 =n-\sum_{j\ge1}\left\lfloor\frac n{2^j}\right\rfloor.
\]
Only finitely many floor terms are nonzero.
\end{proof}

\begin{proposition}[Fractional-part sum]
\label{prop:fractional-weight}
For every \(n\ge0\),
\begin{equation}
 \boxed{\swt(n)=\sum_{j\ge1}\left\{\frac{n}{2^j}\right\}},
 \label{eq:fractional-weight}
\end{equation}
where \(\{x\}=x-\lfloor x\rfloor\).  Hence the apparently real-valued series on the right is
always an integer, and
\begin{equation}
 \eps_n=\cos\!\left(\pi\sum_{j\ge1}\left\{\frac n{2^j}\right\}\right),
 \qquad
 \tau(n)=\frac{1-\cos\!\left(\pi\sum_{j\ge1}\{n/2^j\}\right)}2.
 \label{eq:fractional-cos}
\end{equation}
\end{proposition}

\begin{proof}
The geometric series gives \(\sum_{j\ge1}n/2^j=n\).  Subtract the floor sum in
\eqref{eq:weight-floor}:
\[
 \sum_{j\ge1}\left\{\frac n{2^j}\right\}
 =n-\sum_{j\ge1}\left\lfloor\frac n{2^j}\right\rfloor
 =\swt(n).
\]
The cosine formulas follow because \(\cos(\pi k)=(-1)^k\) for integer \(k\).
\end{proof}

\begin{remark}
Unlike the digit product, \eqref{eq:fractional-weight} is genuinely infinite, although its
tail is a geometric series.  It is a useful bridge between digit sums and periodic
sawtooth functions.
\end{remark}

\subsection{A pure trigonometric sign product}

\begin{theorem}[Sine-sign formula]
\label{thm:sine-sign}
For every \(n\ge0\),
\begin{equation}
 \boxed{
 \eps_n=
 \prod_{j=0}^{\infty}
 \sgn\!\left(\sin\!\frac{\pi(n+\tfrac12)}{2^j}\right).}
 \label{eq:sine-sign}
\end{equation}
Every factor is nonzero, and all sufficiently large factors equal \(+1\), so the product is
well defined without a limiting convention.
\end{theorem}

\begin{proof}
For an integer \(q\) and \(0<\theta<1\),
\(\sgn\sin(\pi(q+\theta))=(-1)^q\).  Moreover
\[
 \left\lfloor\frac{n+1/2}{2^j}\right\rfloor
 =\left\lfloor\frac n{2^j}\right\rfloor
\]
for every \(j\ge0\).  Therefore the product in \eqref{eq:sine-sign} is
\[
 (-1)^{\sum_{j\ge0}\lfloor n/2^j\rfloor}.
\]
By \eqref{eq:weight-floor}, the exponent is
\[
 n+\sum_{j\ge1}\left\lfloor\frac n{2^j}\right\rfloor
 =n+(n-\swt(n))=2n-\swt(n),
\]
which has the same parity as \(\swt(n)\).  Finally, once \(2^j>n+1/2\), the sine argument lies
in \((0,\pi)\), so the factor is \(+1\).
\end{proof}

\section{Factorial valuations, binomial valuations, and carries}
\label{sec:valuations}

\subsection{Legendre's formula as a Thue--Morse formula}

Let \(\vtwo(M)\) denote the exponent of two in a positive integer \(M\).

\begin{proposition}[Factorial-valuation formula]
\label{prop:factorial-valuation}
For every \(n\ge0\),
\begin{equation}
 \vtwo(n!)=\sum_{j\ge1}\left\lfloor\frac n{2^j}\right\rfloor=n-\swt(n).
 \label{eq:legendre2}
\end{equation}
Thus
\begin{equation}
 \boxed{
 \eps_n=(-1)^{n-\vtwo(n!)},
 \qquad
 \tau(n)\equiv n-\vtwo(n!)\pmod2.}
 \label{eq:factorial-tm}
\end{equation}
\end{proposition}

\begin{proof}
Each multiple of \(2\) contributes at least one factor of two to \(n!\), each multiple of
\(4\) contributes one additional factor, and so on, proving the first equality.  The second
is \eqref{eq:weight-floor}; the Thue--Morse formulas follow from the definitions.
\end{proof}

The formula can also be read cumulatively.  Since
\(\vtwo(n!)=\sum_{k=1}^{n}\vtwo(k)\),
\begin{equation}
 \eps_n=(-1)^{\,n+\sum_{k=1}^{n}\vtwo(k)}.
 \label{eq:ruler-cumulative}
\end{equation}
The minus sign in \eqref{eq:factorial-tm} may be changed to a plus sign in an exponent of
\((-1)\).

\subsection{Central binomial coefficients and Catalan numbers}

\begin{theorem}[Central-binomial valuation formula]
\label{thm:central-binomial}
For every \(n\ge0\),
\begin{equation}
 \boxed{\vtwo\binom{2n}{n}=\swt(n)}.
 \label{eq:central-v2}
\end{equation}
Consequently
\begin{equation}
 \boxed{
 \eps_n=(-1)^{\vtwo\binom{2n}{n}},
 \qquad
 \tau(n)\equiv\vtwo\binom{2n}{n}\pmod2.}
 \label{eq:central-tm}
\end{equation}
\end{theorem}

\begin{proof}
Using \eqref{eq:legendre2} and \(\swt(2n)=\swt(n)\),
\[
 \vtwo\binom{2n}{n}
 =\vtwo((2n)!)-2\vtwo(n!)
 =(2n-\swt(n))-2(n-\swt(n))
 =\swt(n).
\]
\end{proof}

\begin{corollary}[Catalan valuation]
\label{cor:catalan}
For \(\Cat_n=\frac1{n+1}\binom{2n}{n}\),
\begin{equation}
 \vtwo(\Cat_n)=\swt(n+1)-1.
 \label{eq:catalan-v2}
\end{equation}
Hence
\begin{equation}
 \boxed{\eps_{n+1}=-(-1)^{\vtwo(\Cat_n)}}.
 \label{eq:catalan-sign}
\end{equation}
\end{corollary}

\begin{proof}
The identity
\begin{equation}
 \swt(n+1)-\swt(n)=1-\vtwo(n+1)
 \label{eq:successor-weight}
\end{equation}
results from changing the trailing one-bits of \(n\) to zero and the next zero-bit to one.
Now
\[
 \vtwo(\Cat_n)=\swt(n)-\vtwo(n+1)=\swt(n+1)-1.
\]
Taking \((-1)\) to both sides gives \eqref{eq:catalan-sign}.
\end{proof}

\subsection{The carry cocycle}

\begin{theorem}[Kummer carry formula at two]
\label{thm:carry}
For \(m,n\ge0\),
\begin{equation}
 \vtwo\binom{m+n}{m}
 =\swt(m)+\swt(n)-\swt(m+n).
 \label{eq:carry-v2}
\end{equation}
This number is exactly the number of carries, counted with propagation, in the binary addition
of \(m\) and \(n\).  In particular,
\begin{equation}
 \boxed{
 \eps_{m+n}=\eps_m\eps_n(-1)^{\vtwo\binom{m+n}{m}}.}
 \label{eq:carry-cocycle}
\end{equation}
Equivalently,
\begin{equation}
 \tau(m+n)\equiv
 \tau(m)+\tau(n)+\vtwo\binom{m+n}{m}\pmod2.
 \label{eq:carry-bit}
\end{equation}
\end{theorem}

\begin{proof}
Legendre's formula gives
\[
 \vtwo\binom{m+n}{m}
 =(m+n-\swt(m+n))-(m-\swt(m))-(n-\swt(n)),
\]
which is \eqref{eq:carry-v2}.  The carry interpretation is the binary case of Kummer's
theorem; it can also be seen directly because each carry reduces the total visible digit sum
by one at its source and may propagate, with the valuation recording the complete reduction.
Reducing \eqref{eq:carry-v2} modulo two and exponentiating \((-1)\) gives
\eqref{eq:carry-cocycle}; the bit form follows from \eqref{eq:conversion}.
\end{proof}

\begin{corollary}[Successor formula from the ruler sequence]
\label{cor:successor-sign}
For every \(n\ge0\),
\begin{equation}
 \boxed{\eps_{n+1}=-(-1)^{\vtwo(n+1)}\eps_n.}
 \label{eq:successor-sign}
\end{equation}
\end{corollary}

\begin{proof}
This is \eqref{eq:successor-weight} exponentiated by \((-1)\).
\end{proof}

\begin{remark}
Equation \eqref{eq:carry-cocycle} says that \(n\mapsto\eps_n\) would be additive-to-
multiplicative were it not for carries.  The parity of the binomial valuation is the exact
correction term.  In cohomological language, carries form a two-cocycle for binary addition.
\end{remark}

\section{Pascal support and an \texorpdfstring{\(r\)}{r}-nomial logarithm family}
\label{sec:pascal}

\subsection{Lucas support as a Boolean lattice}

For nonnegative integers \(k,n\), write \(k\submask n\) when
\(b_j(k)\le b_j(n)\) for every \(j\); in computer language, \(k\) is a bitwise submask of
\(n\).

\begin{theorem}[Lucas parity criterion]
\label{thm:lucas-parity}
For \(0\le k\le n\),
\begin{equation}
 \binom nk\equiv1\pmod2
 \quad\Longleftrightarrow\quad
 k\submask n.
 \label{eq:lucas-submask}
\end{equation}
Therefore the odd entries in row \(n\) are naturally indexed by the subsets of the set of
one-bit positions of \(n\).
\end{theorem}

\begin{proof}
In \(\F_2[z]\),
\[
 (1+z)^n
 =\prod_{j:b_j(n)=1}(1+z)^{2^j}
 =\prod_{j:b_j(n)=1}(1+z^{2^j}).
\]
On expansion, a monomial is obtained by choosing a subset of the set-bit positions, and its
exponent is the integer whose one-bits are precisely that subset.  Thus the nonzero
coefficient positions are exactly the submasks of \(n\).
\end{proof}

\begin{corollary}[Glaisher--Fine count]
\label{cor:glaisher-count}
The number \(O(n)\) of odd binomial coefficients \(\binom nk\) is
\begin{equation}
 O(n)=2^{\swt(n)}.
 \label{eq:glaisher-count}
\end{equation}
\end{corollary}

\begin{proof}
A set with \(\swt(n)\) elements has \(2^{\swt(n)}\) subsets.
\end{proof}

This gives the direct logarithm formula \eqref{eq:pascal-count-direct}; the source formula
\eqref{eq:source-formula} is the same count recovered from a signed sum of the whole row.

\subsection{Odd multinomial coefficients}

Fix an integer \(r\ge2\).  For a weak composition
\(\bm k=(k_1,\ldots,k_r)\) of \(n\), write
\[
 \binom{n}{k_1,\ldots,k_r}=\frac{n!}{k_1!\cdots k_r!}.
\]
Let \(O_r(n)\) be the number of odd multinomial coefficients in the expansion of
\((x_1+\cdots+x_r)^n\).

\begin{theorem}[Multinomial support count]
\label{thm:multinomial-count}
For every \(r\ge2\) and \(n\ge0\),
\begin{equation}
 \boxed{O_r(n)=r^{\swt(n)}.}
 \label{eq:multinomial-count}
\end{equation}
Consequently,
\begin{equation}
 \boxed{
 \eps_n=(-1)^{\log_r O_r(n)},
 \qquad
 \tau(n)=\frac{1-(-1)^{\log_rO_r(n)}}2.}
 \label{eq:multinomial-log-direct}
\end{equation}
The logarithm is exact because \(O_r(n)\) is proved first to be an integral power of \(r\).
\end{theorem}

\begin{proof}
In characteristic two, Frobenius gives
\[
 (x_1+\cdots+x_r)^{2^j}=x_1^{2^j}+\cdots+x_r^{2^j}.
\]
Hence
\begin{equation}
 (x_1+\cdots+x_r)^n
 =\prod_{j:b_j(n)=1}
   \bigl(x_1^{2^j}+\cdots+x_r^{2^j}\bigr)
 \qquad\text{in }\F_2[x_1,\ldots,x_r].
 \label{eq:multinomial-frob-factor}
\end{equation}
For each one-bit \(2^j\) of \(n\), choose one of the \(r\) variables to receive that bit.
Different assignments yield different exponent vectors because binary expansions are unique,
and every nonzero monomial arises this way.  There are \(r\) independent choices for each of
\(\swt(n)\) bits, proving \eqref{eq:multinomial-count}.  Taking the parity of the exponent
\(\swt(n)=\log_rO_r(n)\) proves \eqref{eq:multinomial-log-direct}.
\end{proof}

\subsection{A signed multinomial analogue of section 11.8}

There are \(\binom{n+r-1}{r-1}\) weak compositions of \(n\) into \(r\) parts.  Define
\begin{equation}
 \mathcal X_r(n)=
 \sum_{k_1+\cdots+k_r=n}
 (-1)^{\binom{n}{k_1,\ldots,k_r}}.
 \label{eq:Xr}
\end{equation}

\begin{proposition}[Exact \(r\)-nomial logarithm]
\label{prop:rnomial-exact}
For \(r\ge2\),
\begin{equation}
 \binom{n+r-1}{r-1}-\mathcal X_r(n)
 =2r^{\swt(n)}.
 \label{eq:rnomial-exact}
\end{equation}
Thus
\begin{equation}
 \boxed{
 \tau(n)=
 \frac{1-(-1)^{\displaystyle
  \log_r\!\left(\frac{\binom{n+r-1}{r-1}-\mathcal X_r(n)}2\right)}}2.}
 \label{eq:rnomial-log}
\end{equation}
For \(r=2\), this is equivalent to the binomial--logarithm mechanism of section 11.8.
\end{proposition}

\begin{proof}
A summand of \(\mathcal X_r(n)\) is \(-1\) when the corresponding multinomial coefficient is
odd and \(+1\) when it is even.  If \(T=\binom{n+r-1}{r-1}\) is the total number of summands,
then
\[
 \mathcal X_r(n)=T-2O_r(n).
\]
Apply \eqref{eq:multinomial-count}; the quantity inside the logarithm in
\eqref{eq:rnomial-log} is exactly \(r^{\swt(n)}\).
\end{proof}

\begin{example}
Take \(n=13=1101_2\), so \(\swt(13)=3\).  For \(r=3\), exactly
\(O_3(13)=3^3=27\) coefficients of \((x+y+z)^{13}\) are odd.  The total number of
coefficients is \(\binom{15}{2}=105\), so
\(\mathcal X_3(13)=105-2\cdot27=51\).  The exact logarithm reads
\(\log_3((105-51)/2)=\log_3 27=3\), whose parity gives \(\tau(13)=1\).
\end{example}

\section{Finite and infinite generating products}
\label{sec:products}

\subsection{The finite product polynomial}

Define
\begin{equation}
 P_m(z)=\sum_{n=0}^{2^m-1}\eps_nz^n.
 \label{eq:Pm-def}
\end{equation}

\begin{theorem}[Finite Thue--Morse product]
\label{thm:finite-product}
For every \(m\ge0\),
\begin{equation}
 \boxed{
 P_m(z)=\prod_{j=0}^{m-1}(1-z^{2^j}).}
 \label{eq:finite-product}
\end{equation}
The empty product for \(m=0\) is \(1\).
\end{theorem}

\begin{proof}
The complement-and-concatenate rule gives
\[
 P_{m+1}(z)
 =P_m(z)-z^{2^m}P_m(z)
 =(1-z^{2^m})P_m(z).
\]
Starting from \(P_0(z)=1\) and iterating proves the product.  Equivalently, expanding the
product chooses a subset of the powers \(1,2,\ldots,2^{m-1}\); the resulting exponent is the
integer with that set of one-bits, and the coefficient is \((-1)^{\text{number of choices}}
=\eps_n\).
\end{proof}

\begin{corollary}[Coefficient extraction]
\label{cor:coefficient}
For every \(n\ge0\),
\begin{equation}
 \boxed{
 \eps_n=[z^n]\prod_{j=0}^{\lfloor\log_2n\rfloor}(1-z^{2^j})}
 \qquad(n\ge1),
 \label{eq:coefficient-formula}
\end{equation}
with \(\eps_0=[z^0]1=1\).
\end{corollary}

\subsection{The infinite product and the bit generating function}

Let
\begin{equation}
 E(z)=\sum_{n\ge0}\eps_nz^n,
 \qquad
 T(z)=\sum_{n\ge0}\tau(n)z^n.
 \label{eq:E-T-def}
\end{equation}

\begin{theorem}[Infinite product]
\label{thm:infinite-product}
For \(|z|<1\),
\begin{equation}
 \boxed{
 E(z)=\prod_{j=0}^{\infty}(1-z^{2^j}).}
 \label{eq:infinite-product}
\end{equation}
The product converges uniformly on compact subsets of the unit disk.  Moreover,
\begin{equation}
 \boxed{
 T(z)=\frac{1}{2(1-z)}-\frac12\prod_{j=0}^{\infty}(1-z^{2^j}).}
 \label{eq:T-product}
\end{equation}
\end{theorem}

\begin{proof}
For \(|z|\le\rho<1\), the series \(\sum_j|z|^{2^j}\) converges, so the product converges
uniformly on the compact disk.  Every coefficient stabilizes in the finite products
\(P_m\), and \cref{thm:finite-product} identifies it as \(\eps_n\).  Finally,
\(\tau(n)=(1-\eps_n)/2\), so
\[
 T(z)=\frac12\sum_{n\ge0}z^n-\frac12E(z).
\]
\end{proof}

\subsection{Mahler functional equations and an iterated quotient}

\begin{proposition}[Two Mahler equations]
\label{prop:mahler}
For \(|z|<1\),
\begin{align}
 E(z)&=(1-z)E(z^2),                                      \label{eq:E-mahler}\\
 T(z)&=(1-z)T(z^2)+\frac{z}{1-z^2}.                      \label{eq:T-mahler}
\end{align}
\end{proposition}

\begin{proof}
Separate the first factor in \eqref{eq:infinite-product} to obtain
\eqref{eq:E-mahler}.  Substitute \eqref{eq:T-product}, or split the even and odd
coefficients using \eqref{eq:tau-evenodd}, to obtain \eqref{eq:T-mahler}.
\end{proof}

Write \(P_0(z)=1\) and \(P_j(z)=\prod_{h=0}^{j-1}(1-z^{2^h})\) for \(j\ge1\).

\begin{theorem}[Lacunary quotient expansion]
\label{thm:lacunary-quotient}
For \(|z|<1\),
\begin{equation}
 \boxed{
 T(z)=\sum_{j=0}^{\infty}
 P_j(z)\frac{z^{2^j}}{1-z^{2^{j+1}}}.}
 \label{eq:lacunary-quotient}
\end{equation}
\end{theorem}

\begin{proof}
Iterating \eqref{eq:T-mahler} \(M\) times gives
\[
 T(z)=\sum_{j=0}^{M-1}P_j(z)\frac{z^{2^j}}{1-z^{2^{j+1}}}
      +P_M(z)T(z^{2^M}).
\]
On a compact subdisk, \(P_M(z)\) remains bounded while \(T(z^{2^M})\to T(0)=0\), so the
remainder tends to zero.
\end{proof}

\subsection{The Thue--Morse constant}

The binary real whose digits are the Thue--Morse bits is
\begin{equation}
 \kappa=0.\tau(0)\tau(1)\tau(2)\cdots{}_2
       =\sum_{n\ge0}\frac{\tau(n)}{2^{n+1}}.
 \label{eq:kappa-def}
\end{equation}

\begin{corollary}[A lacunary product for \(\kappa\)]
\label{cor:kappa}
\begin{equation}
 \boxed{
 \kappa=\frac{2-\displaystyle\prod_{j=0}^{\infty}
                    \left(1-2^{-2^j}\right)}4.}
 \label{eq:kappa-product}
\end{equation}
\end{corollary}

\begin{proof}
By definition, \(\kappa=\tfrac12T(1/2)\).  Substitute \eqref{eq:T-product}.
\end{proof}

\section{Ordinary Fourier and Walsh formulas}
\label{sec:fourier}

\subsection{Trigonometric factorization on the unit circle}

\begin{proposition}[Finite Fourier product]
\label{prop:unit-circle}
For real \(x\),
\begin{align}
 \sum_{n=0}^{2^m-1}\eps_n\e^{2\pi\ii nx}
 &=\prod_{j=0}^{m-1}\left(1-\e^{2\pi\ii2^jx}\right),
 \label{eq:fourier-product}\\
 &=(-2\ii)^m\e^{\pi\ii(2^m-1)x}
   \prod_{j=0}^{m-1}\sin(\pi2^jx).                         \label{eq:sine-fourier}
\end{align}
In particular,
\begin{equation}
 \left|\sum_{n=0}^{2^m-1}\eps_n\e^{2\pi\ii nx}\right|
 =2^m\prod_{j=0}^{m-1}|\sin(\pi2^jx)|.
 \label{eq:fourier-magnitude}
\end{equation}
\end{proposition}

\begin{proof}
Substitute \(z=\e^{2\pi\ii x}\) in \eqref{eq:finite-product} and use
\(1-\e^{2\ii y}=-2\ii\e^{\ii y}\sin y\).  The phase exponents sum to
\(\sum_{j=0}^{m-1}2^j=2^m-1\).
\end{proof}

\subsection{A discrete Fourier inversion formula for each bit}

Let \(\omega_m=\e^{2\pi\ii/2^m}\).

\begin{theorem}[Root-of-unity formula]
\label{thm:dft}
For \(0\le n<2^m\),
\begin{equation}
 \boxed{
 \eps_n=\frac1{2^m}
 \sum_{r=0}^{2^m-1}
 \left(\prod_{j=0}^{m-1}(1-\omega_m^{r2^j})\right)
 \omega_m^{-rn}.}
 \label{eq:dft-inversion}
\end{equation}
All terms with even \(r\) vanish, so the sum may be restricted to odd \(r\).
\end{theorem}

\begin{proof}
This is the inverse discrete Fourier transform of the coefficient vector of \(P_m\).  If
\(r\) is even, the factor with \(j=m-1\) is
\(1-\omega_m^{r2^{m-1}}=1-(-1)^r=0\).
\end{proof}

\begin{remark}
Formula \eqref{eq:dft-inversion} is deliberately extravagant as an evaluator: it replaces one
parity computation by a complex Fourier sum.  Its value is conceptual.  It exhibits the bit
as a Fourier coefficient of a completely factored cyclotomic object.
\end{remark}

\subsection{Walsh--Hadamard concentration}

For \(\bm a,\bm b\in\F_2^m\), write
\(\chi_{\bm a}(\bm b)=(-1)^{\bm a\cdot\bm b}\).

\begin{theorem}[A single Walsh frequency]
\label{thm:walsh}
Let \(\bm1=(1,\ldots,1)\in\F_2^m\).  Then
\begin{equation}
 \eps_n=\chi_{\bm1}(\bm b(n)),
 \label{eq:walsh-character}
\end{equation}
and the Walsh transform is
\begin{equation}
 \sum_{n=0}^{2^m-1}\eps_n\chi_{\bm a}(\bm b(n))
 =\begin{cases}
   2^m,&\bm a=\bm1,\\
   0,&\bm a\ne\bm1.
  \end{cases}
 \label{eq:walsh-transform}
\end{equation}
\end{theorem}

\begin{proof}
The first formula is the definition of parity.  The summand in the transform is
\(\chi_{\bm a+\bm1}(\bm b(n))\).  The sum of a nontrivial character over the finite group
\(\F_2^m\) is zero, while the trivial character sums to \(2^m\).
\end{proof}

\section{Prouhet cancellation, finite differences, and exact moments}
\label{sec:prouhet}

\subsection{The zero of order \(m\) at one}

\begin{lemma}[Exact vanishing order]
\label{lem:order-m}
The polynomial \(P_m(z)\) has a zero of exact order \(m\) at \(z=1\), and
\begin{equation}
 P_m(z)=(-1)^m2^{\binom m2}(z-1)^m+O((z-1)^{m+1}).
 \label{eq:Pm-local}
\end{equation}
\end{lemma}

\begin{proof}
Every factor \(1-z^{2^j}\) has a simple zero at one and local expansion
\(-2^j(z-1)+O((z-1)^2)\).  Multiply the \(m\) leading terms and use
\(\sum_{j=0}^{m-1}j=\binom m2\).
\end{proof}

\begin{theorem}[Prouhet cancellation]
\label{thm:prouhet}
For every polynomial \(p\in\C[x]\) of degree less than \(m\),
\begin{equation}
 \boxed{
 \sum_{n=0}^{2^m-1}\eps_np(n)=0.}
 \label{eq:prouhet-poly}
\end{equation}
Equivalently, if
\[
 A_0=\{0\le n<2^m:\tau(n)=0\},
 \qquad
 A_1=\{0\le n<2^m:\tau(n)=1\},
\]
then for \(0\le r<m\),
\begin{equation}
 \sum_{n\in A_0}n^r=\sum_{n\in A_1}n^r.
 \label{eq:prouhet-partition}
\end{equation}
\end{theorem}

\begin{proof}
For \(r\ge0\), differentiating \eqref{eq:Pm-def} gives
\[
 P_m^{(r)}(1)=\sum_{n<2^m}\eps_n\fall nr.
\]
By \cref{lem:order-m}, this vanishes for \(r<m\).  The falling factorials
\(1,x,\fall x2,\ldots\) form a basis of the polynomial ring, proving
\eqref{eq:prouhet-poly}.  Taking \(p(x)=x^r\) and separating the signs gives
\eqref{eq:prouhet-partition}.
\end{proof}

\begin{corollary}[Affine Prouhet identities]
\label{cor:affine-prouhet}
For \(a,h\in\C\) and \(\deg p<m\),
\begin{equation}
 \sum_{n<2^m}\eps_np(a+hn)=0.
 \label{eq:affine-prouhet}
\end{equation}
\end{corollary}

\begin{proof}
The function \(n\mapsto p(a+hn)\) is a polynomial in \(n\) of degree at most \(\deg p\).
\end{proof}

\subsection{The first nonzero moments}

\begin{theorem}[First surviving power and binomial moments]
\label{thm:first-moment}
For \(m\ge0\),
\begin{align}
 \sum_{n=0}^{2^m-1}\eps_n n^m
 &=(-1)^m m!\,2^{\binom m2},                              \label{eq:first-power-moment}\\
 \sum_{n=0}^{2^m-1}\eps_n\binom nm
 &=(-1)^m2^{\binom m2}.                                   \label{eq:first-binomial-moment}
\end{align}
\end{theorem}

\begin{proof}
By \eqref{eq:Pm-local},
\(P_m^{(m)}(1)=(-1)^m m!2^{\binom m2}\).  This derivative is the sum of
\(\eps_n\fall nm\), and \(n^m-\fall nm\) has degree less than \(m\), so it vanishes under the
Prouhet sum.  This proves \eqref{eq:first-power-moment}.  Divide the falling-factorial identity
by \(m!\) to obtain \eqref{eq:first-binomial-moment}.
\end{proof}

\subsection{A finite-difference operator identity}

For \(a\in\C\), let \(T_a f(x)=f(x+a)\).

\begin{theorem}[Dyadic finite-difference factorization]
\label{thm:operator}
For every function \(f\) on which the shifts are defined,
\begin{equation}
 \boxed{
 \sum_{n=0}^{2^m-1}\eps_n f(x+nh)
 =\prod_{j=0}^{m-1}\bigl(I-T_{2^jh}\bigr)f(x).}
 \label{eq:operator-factor}
\end{equation}
\end{theorem}

\begin{proof}
Expand the product of operators.  A subset \(S\subseteq\{0,\ldots,m-1\}\) contributes the
shift by \(h\sum_{j\in S}2^j\) with sign \((-1)^{|S|}\).  The subset sums are precisely the
integers \(0\le n<2^m\), and \((-1)^{|S|}=\eps_n\).
\end{proof}

\begin{corollary}[Exact multiple-integral remainder]
\label{cor:integral-remainder}
Suppose \(h\ge0\) and \(f\) is \(m\) times continuously differentiable on an interval
containing all relevant points.  Then
\begin{align}
 \sum_{n<2^m}\eps_nf(x+nh)
 &=(-1)^m
 \int_0^h\int_0^{2h}\cdots\int_0^{2^{m-1}h}
 f^{(m)}(x+t_0+\cdots+t_{m-1})
 \dd t_{m-1}\cdots\dd t_0.                                \label{eq:integral-remainder}
\end{align}
In particular,
\begin{equation}
 \left|\sum_{n<2^m}\eps_nf(x+nh)\right|
 \le 2^{\binom m2}h^m
 \sup_{x\le y\le x+(2^m-1)h}|f^{(m)}(y)|.
 \label{eq:difference-bound}
\end{equation}
\end{corollary}

\begin{proof}
Use \((I-T_a)f(x)=-\int_0^af'(x+t)\dd t\) successively in
\eqref{eq:operator-factor}.  The volume of the integration box is
\(h\cdot2h\cdots2^{m-1}h=2^{\binom m2}h^m\).
\end{proof}

\subsection{All higher binomial moments}

Define
\begin{equation}
 B_{m,r}=\sum_{n=0}^{2^m-1}\eps_n\binom nr.
 \label{eq:Bmr}
\end{equation}

\begin{theorem}[Coefficient formula for every binomial moment]
\label{thm:all-binomial-moments}
For \(\ell\ge0\),
\begin{equation}
 \boxed{
 B_{m,m+\ell}
 =(-1)^m2^{\binom m2}[u^\ell]
 \prod_{j=0}^{m-1}
 \frac{(1+u)^{2^j}-1}{2^ju}.}
 \label{eq:all-binomial-moments}
\end{equation}
For \(r<m\), \(B_{m,r}=0\).
\end{theorem}

\begin{proof}
Expanding \(P_m(1+u)\) in two ways gives
\[
 P_m(1+u)
 =\sum_{n<2^m}\eps_n\sum_{r\ge0}\binom nru^r
 =\sum_{r\ge0}B_{m,r}u^r,
\]
while the product factorization gives
\[
 P_m(1+u)
 =\prod_{j=0}^{m-1}\bigl(1-(1+u)^{2^j}\bigr)
 =(-1)^m2^{\binom m2}u^m
  \prod_{j=0}^{m-1}\frac{(1+u)^{2^j}-1}{2^ju}.
\]
Compare coefficients.
\end{proof}

\subsection{All higher power moments in closed coefficient form}

Let \(B_{2r}\) be the Bernoulli numbers with \(B_2=1/6\), and put
\begin{equation}
 A_m=\frac{2^m-1}{2},
 \qquad
 c_{m,r}=\frac{B_{2r}}{2r(2r)!}\,
          \frac{2^{2rm}-1}{2^{2r}-1}\quad(r\ge1).
 \label{eq:Am-cmr}
\end{equation}

\begin{theorem}[Exponential moment factorization]
\label{thm:exponential-moment}
As a formal power series in \(t\), and analytically near \(t=0\),
\begin{equation}
 \boxed{
 \sum_{n=0}^{2^m-1}\eps_n\e^{nt}
 =(-1)^m2^{\binom m2}t^m
  \exp\!\left(A_mt+\sum_{r\ge1}c_{m,r}t^{2r}\right).}
 \label{eq:exponential-factor}
\end{equation}
Consequently, for every \(\ell\ge0\),
\begin{equation}
 \boxed{
 \sum_{n=0}^{2^m-1}\eps_n n^{m+\ell}
 =(-1)^m2^{\binom m2}(m+\ell)!
 [t^\ell]\exp\!\left(A_mt+\sum_{r\ge1}c_{m,r}t^{2r}\right).}
 \label{eq:all-power-moments}
\end{equation}
\end{theorem}

\begin{proof}
Substitute \(z=\e^t\) in \eqref{eq:finite-product}.  For \(a>0\),
\[
 1-\e^{at}=-at\,\e^{at/2}\frac{\sinh(at/2)}{at/2}.
\]
Multiplying over \(a=2^0,\ldots,2^{m-1}\) produces the prefactor and the linear exponential
\(A_mt\).  The standard Bernoulli expansion
\[
 \log\frac{\sinh x}{x}
 =\sum_{r\ge1}\frac{2^{2r-1}B_{2r}}{r(2r)!}x^{2r}
\]
then yields \(c_{m,r}\), because
\(\sum_{j=0}^{m-1}2^{2rj}=(2^{2rm}-1)/(2^{2r}-1)\).  Finally compare the coefficient of
\(t^{m+\ell}\) with
\(\sum_n\eps_n\e^{nt}=\sum_q(\sum_n\eps_nn^q)t^q/q!\).
\end{proof}

The first three surviving moments are therefore
\begin{align}
 \sum_{n<2^m}\eps_nn^m
 &=(-1)^m2^{\binom m2}m!,                                  \label{eq:mom0}\\
 \sum_{n<2^m}\eps_nn^{m+1}
 &=(-1)^m2^{\binom m2}(m+1)!\frac{2^m-1}{2},               \label{eq:mom1}\\
 \sum_{n<2^m}\eps_nn^{m+2}
 &=(-1)^m2^{\binom m2}(m+2)!
   \left(\frac{(2^m-1)^2}{8}+\frac{4^m-1}{72}\right).     \label{eq:mom2}
\end{align}

\section{Sparse Prouhet identities on an arbitrary Pascal row}
\label{sec:sparse}

The ordinary Prouhet identity uses all integers below \(2^m\), which are the submasks of
\(2^m-1\).  Lucas's criterion suggests replacing \(2^m-1\) by an arbitrary \(n\).

Let
\begin{equation}
 J(n)=\{j\ge0:b_j(n)=1\},
 \qquad
 s=|J(n)|=\swt(n),
 \qquad
 \beta(n)=\sum_{j\in J(n)}j.
 \label{eq:J-beta}
\end{equation}

\begin{theorem}[Pascal-support product]
\label{thm:submask-product}
For every \(n\ge0\),
\begin{equation}
 \boxed{
 R_n(z):=\sum_{k\submask n}\eps_kz^k
 =\prod_{j\in J(n)}(1-z^{2^j}).}
 \label{eq:submask-product}
\end{equation}
Equivalently, by Lucas's criterion,
\begin{equation}
 R_n(z)=\sum_{k=0}^{n}
 \eps_k\,\1_{\binom nk\text{ odd}}\,z^k.
 \label{eq:pascal-support-polynomial}
\end{equation}
\end{theorem}

\begin{proof}
Expanding the product selects a subset \(S\subseteq J(n)\).  It contributes the exponent
\(k=\sum_{j\in S}2^j\), which is exactly a submask of \(n\), with coefficient
\((-1)^{|S|}=\eps_k\).  Equation \eqref{eq:pascal-support-polynomial} follows from
\cref{thm:lucas-parity}.
\end{proof}

\begin{corollary}[Signed balance on odd Pascal positions]
\label{cor:pascal-balance}
For \(n>0\),
\begin{equation}
 \boxed{
 \sum_{\substack{0\le k\le n\\ \binom nk\text{ odd}}}\eps_k=0.}
 \label{eq:pascal-balance}
\end{equation}
Thus among the odd positions in every nonzero Pascal row, exactly half have \(\tau(k)=0\) and
half have \(\tau(k)=1\).
\end{corollary}

\begin{proof}
Set \(z=1\) in \eqref{eq:submask-product}.  Because \(n>0\), the product has at least one
factor \(1-1=0\).  Since there are \(2^{\swt(n)}\) odd positions, signed balance means the two
classes each have size \(2^{\swt(n)-1}\).
\end{proof}

\begin{theorem}[Sparse Prouhet theorem]
\label{thm:sparse-prouhet}
Let \(n>0\) and \(s=\swt(n)\).  For every polynomial \(p\) of degree less than \(s\),
\begin{equation}
 \boxed{
 \sum_{k\submask n}\eps_kp(k)=0.}
 \label{eq:sparse-prouhet}
\end{equation}
The first surviving power moment is
\begin{equation}
 \boxed{
 \sum_{k\submask n}\eps_kk^s=(-1)^s s!\,2^{\beta(n)}.}
 \label{eq:sparse-first-moment}
\end{equation}
\end{theorem}

\begin{proof}
Each factor of \(R_n(z)\) has a simple zero at one, so \(R_n\) has a zero of order
\(s\).  Its leading local term is
\[
 R_n(z)=(-1)^s2^{\beta(n)}(z-1)^s+O((z-1)^{s+1}).
\]
The derivative argument from \cref{thm:prouhet,thm:first-moment} now applies verbatim, with
the sum restricted to submasks of \(n\).
\end{proof}

\begin{remark}[Boolean-lattice Möbius inversion]
The submasks of \(n\) form a Boolean lattice on the ground set \(J(n)\).  The coefficient
\(\eps_k=(-1)^{\swt(k)}\) is its Möbius function from the bottom element to the subset
corresponding to \(k\).  Equation \eqref{eq:sparse-prouhet} is therefore a weighted
finite-difference identity on a nonuniform grid with step sizes \(2^j\), \(j\in J(n)\).
\end{remark}

\begin{corollary}[Operator form on a Pascal support]
\label{cor:sparse-operator}
For every function \(f\),
\begin{equation}
 \sum_{k\submask n}\eps_k f(x+hk)
 =\prod_{j\in J(n)}(I-T_{2^jh})f(x).
 \label{eq:sparse-operator}
\end{equation}
If \(f\) is \(s\) times continuously differentiable and \(h\ge0\), then the absolute value is
at most
\begin{equation}
 2^{\beta(n)}h^s
 \sup_{x\le y\le x+nh}|f^{(s)}(y)|.
 \label{eq:sparse-bound}
\end{equation}
\end{corollary}

\begin{proof}
Expand the product as in \cref{thm:operator}; then apply the fundamental theorem of calculus
once for each set bit of \(n\).
\end{proof}

\begin{example}
Let \(n=13=1101_2\).  Then \(J(13)=\{0,2,3\}\), \(s=3\), and \(\beta(13)=5\).  The odd
positions in row \(13\) are the submasks
\[
 0,1,4,5,8,9,12,13.
\]
Their signed sums against \(1,k,k^2\) vanish, while
\[
 \sum_{k\submask13}\eps_kk^3=(-1)^3\,3!\,2^5=-192.
\]
This is a Prouhet partition on a sparse, nonconsecutive set dictated by one Pascal row.
\end{example}

\section{Exact prefix sums, balance, and dyadic blocks}
\label{sec:prefix}

Define the signed prefix discrepancy and the number of one-bits in the sequence prefix by
\begin{equation}
 S(N)=\sum_{n=0}^{N-1}\eps_n,
 \qquad
 A(N)=\sum_{n=0}^{N-1}\tau(n).
 \label{eq:S-A}
\end{equation}

\begin{theorem}[Exact prefix discrepancy]
\label{thm:prefix}
For every \(q\ge0\),
\begin{equation}
 \boxed{
 S(2q)=0,
 \qquad
 S(2q+1)=\eps_q.}
 \label{eq:prefix-exact}
\end{equation}
Consequently
\begin{equation}
 \boxed{
 A(N)=\frac{N-S(N)}2
 =\begin{cases}
    N/2,&N\text{ even},\\[0.25em]
    \bigl(N-\eps_{(N-1)/2}\bigr)/2,&N\text{ odd}.
  \end{cases}}
 \label{eq:ones-prefix}
\end{equation}
\end{theorem}

\begin{proof}
Every adjacent pair cancels:
\(\eps_{2n}+\eps_{2n+1}=\eps_n-\eps_n=0\).  An odd-length prefix has the additional term
\(\eps_{2q}=\eps_q\).  The formula for \(A\) follows from
\(\eps_n=1-2\tau(n)\).
\end{proof}

\begin{corollary}[Interval discrepancy]
\label{cor:interval-discrepancy}
For integers \(0\le a<b\),
\begin{equation}
 \sum_{n=a}^{b-1}\eps_n=S(b)-S(a),
 \qquad
 \left|\sum_{n=a}^{b-1}\eps_n\right|\le2.
 \label{eq:interval-discrepancy}
\end{equation}
Thus every finite interval contains numbers of even and odd binary weight in counts differing
by at most two.
\end{corollary}

\begin{corollary}[Balanced aligned dyadic blocks]
\label{cor:dyadic-block-balance}
For \(m\ge1\), \(a\ge0\),
\begin{equation}
 \sum_{r=0}^{2^m-1}\eps_{2^ma+r}=0.
 \label{eq:dyadic-block-balance}
\end{equation}
More generally, for every \(r<2^m\),
\begin{equation}
 \eps_{2^mn+r}=\eps_r\eps_n.
 \label{eq:power-two-progressions}
\end{equation}
\end{corollary}

\begin{proof}
Apply \cref{thm:concat}.  The block sum is
\(\eps_a\sum_{r<2^m}\eps_r=\eps_aP_m(1)=0\).
\end{proof}

\section{Autocorrelation recurrences and finite Riesz products}
\label{sec:correlation}

\subsection{Exact dyadic average recurrences}

For \(m,k\ge0\), define the dyadic correlation average
\begin{equation}
 \eta_m(k)=\frac1{2^m}\sum_{n=0}^{2^m-1}\eps_n\eps_{n+k}.
 \label{eq:eta-m}
\end{equation}
The second index is not reduced modulo \(2^m\); the infinite sequence is used.

\begin{proposition}[Finite correlation recursion]
\label{prop:finite-correlation}
For \(m\ge0\) and \(r\ge0\),
\begin{align}
 \eta_{m+1}(2r)&=\eta_m(r),                                  \label{eq:eta-even}\\
 \eta_{m+1}(2r+1)&=-\frac{\eta_m(r)+\eta_m(r+1)}2.           \label{eq:eta-odd}
\end{align}
\end{proposition}

\begin{proof}
Split the defining sum into even and odd \(n\).  For an even shift,
\[
 \eps_{2j}\eps_{2j+2r}=\eps_j\eps_{j+r},
 \qquad
 \eps_{2j+1}\eps_{2j+1+2r}=(-\eps_j)(-\eps_{j+r}),
\]
so the two halves agree.  For an odd shift,
\[
 \eps_{2j}\eps_{2j+2r+1}=-\eps_j\eps_{j+r},
 \qquad
 \eps_{2j+1}\eps_{2j+2r+2}=-\eps_j\eps_{j+r+1}.
\]
Average the two halves.
\end{proof}

\begin{theorem}[Limiting autocorrelation]
\label{thm:limiting-correlation}
For every fixed \(k\ge0\), the limit
\begin{equation}
 \eta(k)=\lim_{m\to\infty}\eta_m(k)
 \label{eq:eta-limit}
\end{equation}
exists.  It is characterized by
\begin{align}
 \eta(0)&=1,                                                  \label{eq:eta0}\\
 \eta(2r)&=\eta(r),                                           \label{eq:eta-limit-even}\\
 \eta(2r+1)&=-\frac{\eta(r)+\eta(r+1)}2.                     \label{eq:eta-limit-odd}
\end{align}
In particular,
\begin{equation}
 \eta(1)=-\frac13,
 \qquad
 \eta(2)=-\frac13,
 \qquad
 \eta(3)=\frac13,
 \qquad
 \eta(5)=0.
 \label{eq:eta-examples}
\end{equation}
\end{theorem}

\begin{proof}
We induct on \(k\).  The case \(k=0\) is constant.  For \(k=1\),
\eqref{eq:eta-odd} gives
\[
 \eta_{m+1}(1)=-\frac{1+\eta_m(1)}2,
\]
a contraction with fixed point \(-1/3\).  For \(k>1\), if \(k=2r\), convergence follows from
that of the smaller index \(r\).  If \(k=2r+1\), then both \(r\) and \(r+1\) are smaller than
\(k\), so convergence follows from \eqref{eq:eta-odd} and the induction hypothesis.  Passing
to the limit proves the recurrences and the sample values.
\end{proof}

\subsection{The finite Riesz product}

Define the nonnegative trigonometric polynomial
\begin{equation}
 g_m(x)=2^{-m}\left|P_m(\e^{2\pi\ii x})\right|^2.
 \label{eq:gm-def}
\end{equation}

\begin{theorem}[Riesz-product density]
\label{thm:riesz-finite}
For every \(m\ge0\),
\begin{equation}
 \boxed{
 g_m(x)=\prod_{j=0}^{m-1}\bigl(1-\cos(2\pi2^jx)\bigr).}
 \label{eq:riesz-product}
\end{equation}
Moreover, \(g_m(x)\dd x\) is a probability measure on \([0,1]\).
\end{theorem}

\begin{proof}
Because \(|1-\e^{\ii\theta}|^2=2(1-\cos\theta)\), the finite product
\eqref{eq:finite-product} gives \eqref{eq:riesz-product} after division by \(2^m\).  Parseval's
identity yields
\[
 \int_0^1g_m(x)\dd x
 =2^{-m}\sum_{n=0}^{2^m-1}|\eps_n|^2=1.
\]
\end{proof}

The weak limit of these probability measures is the classical Thue--Morse spectral (or
diffraction) measure.  Its Fourier coefficients are the autocorrelations \(\eta(k)\); the
boundary difference between the coefficient of \(g_m\) and \(\eta_m(k)\) is
\(O(k/2^m)\).  The deeper fact that the limiting measure is purely singular continuous is a
literature theorem; see Baake and Grimm \cite{baakegrimm2009} and Queff\'elec
\cite{queffelec2010}.

\section{The bit series over \texorpdfstring{\(\F_2\)}{F2}: an Artin--Schreier formula}
\label{sec:F2}

The ordinary signed series becomes trivial modulo two because \(-1=1\).  The zero--one bit
series retains the information.  Work in the formal power-series ring \(\F_2[[x]]\) and put
\begin{equation}
 \Theta(x)=\sum_{n\ge0}\tau(n)x^n.
 \label{eq:Theta}
\end{equation}

\begin{proposition}[Characteristic-two functional equation]
\label{prop:F2-functional}
In \(\F_2[[x]]\),
\begin{equation}
 \Theta(x)=(1+x)\Theta(x^2)+\frac{x}{1+x^2}.
 \label{eq:F2-functional}
\end{equation}
\end{proposition}

\begin{proof}
Split into even and odd coefficients and use
\(\tau(2n)=\tau(n)\), \(\tau(2n+1)=1+\tau(n)\) in characteristic two:
\[
 \Theta(x)=\Theta(x^2)+x\sum_{n\ge0}(1+\tau(n))x^{2n}
 =(1+x)\Theta(x^2)+\frac{x}{1+x^2}.
\]
\end{proof}

\begin{theorem}[Quadratic algebraic equation]
\label{thm:F2-quadratic}
The bit series satisfies
\begin{equation}
 \boxed{
 (1+x)^3\Theta(x)^2+(1+x)^2\Theta(x)+x=0.}
 \label{eq:F2-quadratic}
\end{equation}
Equivalently, with \(Y(x)=(1+x)\Theta(x)\),
\begin{equation}
 \boxed{
 Y(x)^2+Y(x)=\frac{x}{1+x}.}
 \label{eq:artin-schreier}
\end{equation}
\end{theorem}

\begin{proof}
Frobenius gives \(\Theta(x^2)=\Theta(x)^2\) in characteristic two.  Substitute this into
\eqref{eq:F2-functional} and use \(1+x^2=(1+x)^2\); clearing the denominator yields
\eqref{eq:F2-quadratic}.  Multiplying out \(Y=(1+x)\Theta\) gives
\eqref{eq:artin-schreier}.
\end{proof}

\begin{theorem}[Explicit Artin--Schreier sum]
\label{thm:F2-explicit}
The unique solution of \eqref{eq:artin-schreier} with zero constant term is
\begin{equation}
 Y(x)=\sum_{j=0}^{\infty}\left(\frac{x}{1+x}\right)^{2^j}
     =\sum_{j=0}^{\infty}\frac{x^{2^j}}{1+x^{2^j}}.
 \label{eq:Y-explicit}
\end{equation}
Therefore
\begin{equation}
 \boxed{
 \Theta(x)=\frac1{1+x}
 \sum_{j=0}^{\infty}\frac{x^{2^j}}{1+x^{2^j}}}
 \qquad\text{in }\F_2[[x]].
 \label{eq:Theta-explicit}
\end{equation}
\end{theorem}

\begin{proof}
Let \(A=x/(1+x)\).  The series \(\sum_{j\ge0}A^{2^j}\) is well defined in the
\(x\)-adic topology because the lowest degree doubles at every step.  Squaring shifts the
index, so
\[
 \left(\sum_{j\ge0}A^{2^j}\right)^2+
 \sum_{j\ge0}A^{2^j}=A.
\]
Also \(A^{2^j}=x^{2^j}/(1+x)^{2^j}=x^{2^j}/(1+x^{2^j})\).  If two zero-constant solutions
existed, their difference \(Z\) would satisfy \(Z^2+Z=Z(Z+1)=0\).  Since
\(\F_2[[x]]\) is an integral domain and a zero-constant series cannot equal one, \(Z=0\).
Divide by \(1+x\) to obtain \eqref{eq:Theta-explicit}.
\end{proof}

\begin{remark}
Christol's theorem characterizes automatic sequences over a finite field by algebraicity of
their generating series.  Here no appeal to the theorem is needed: the quadratic equation is
obtained directly.  Conversely, \eqref{eq:F2-quadratic} is an especially small algebraic
certificate of the two-state automaton.
\end{remark}

\section{Evaluated products and other analytic constants}
\label{sec:constants}

The generating product \eqref{eq:infinite-product} is one source of constants.  A different
family uses the signed sequence as exponents of rational functions of \(n\).

\begin{theorem}[Woods--Robbins product; literature formula]
\label{thm:woods-robbins}
With products interpreted as limits of their natural partial products,
\begin{equation}
 \boxed{
 \prod_{n=0}^{\infty}
 \left(\frac{2n+1}{2n+2}\right)^{\eps_n}
 =\frac1{\sqrt2}.}
 \label{eq:woods-robbins}
\end{equation}
\end{theorem}

A convenient way to organize the proof is to introduce a two-parameter product
\[
 F(a,b)=\prod_{n\ge0}\left(\frac{n+a}{n+b}\right)^{\eps_n}
\]
for admissible positive parameters and split its factors into even and odd indices.  The basic
recurrence is
\begin{equation}
 F(a,b)=F\!\left(\frac a2,\frac b2\right)
        F\!\left(\frac{b+1}{2},\frac{a+1}{2}\right),
 \label{eq:F-functional-product}
\end{equation}
from which specializations and cancellations yield \eqref{eq:woods-robbins}.  The convergence
and a broad collection of further evaluations are developed by Allouche, Riasat, and Shallit
\cite{alloucheriasatshallit2019}.  We label this result as literature rather than reproduce
their full product calculus.

\begin{remark}
The product \eqref{eq:woods-robbins} can also be recognized through dyadic partial products:
\[
 \prod_{n=0}^{2^m-1}\left(\frac{2n+1}{2n+2}\right)^{\eps_n}
 =\prod_{k=0}^{2^{m+1}-1}(k+1)^{\eps_k}.
\]
This identity follows by pairing the even and odd indices on the right.  Thus the constant is
a renormalized factorial-like limit controlled entirely by the recursion
\(\eps_{2n}=\eps_n\), \(\eps_{2n+1}=-\eps_n\).
\end{remark}

\section{A base-\texorpdfstring{\(q\)}{q} extension}
\label{sec:baseq}

The binary sequence is the first member of a root-of-unity family.  This section is included
both for perspective and because it shows which product and Prouhet arguments are genuinely
binary and which are not.

Let \(q\ge2\), let \(s_q(n)\) be the sum of the base-\(q\) digits of \(n\), and let
\(\zeta\ne1\) be a \(q\)-th root of unity.  Define
\begin{equation}
 u_n=\zeta^{s_q(n)}.
 \label{eq:generalized-TM}
\end{equation}

\begin{theorem}[Generalized digit recursion and product]
\label{thm:baseq-product}
For \(0\le r<q\),
\begin{equation}
 u_{qn+r}=\zeta^r u_n.
 \label{eq:baseq-recursion}
\end{equation}
Moreover,
\begin{equation}
 \boxed{
 \sum_{n=0}^{q^m-1}u_nz^n
 =\prod_{j=0}^{m-1}
   \left(\sum_{r=0}^{q-1}\zeta^r z^{rq^j}\right).}
 \label{eq:baseq-product}
\end{equation}
\end{theorem}

\begin{proof}
Appending the base-\(q\) digit \(r\) increases the digit sum by \(r\), proving the recurrence.
Expanding the product chooses one digit \(r\) independently at each position \(j\); the
exponent is the corresponding base-\(q\) integer, and the coefficient is \(\zeta\) raised to
its digit sum.
\end{proof}

\begin{corollary}[Generalized Prouhet cancellation]
\label{cor:baseq-prouhet}
If \(\zeta\ne1\) is a \(q\)-th root of unity, then for every polynomial \(p\) with
\(\deg p<m\),
\begin{equation}
 \sum_{n=0}^{q^m-1}\zeta^{s_q(n)}p(n)=0.
 \label{eq:baseq-prouhet}
\end{equation}
\end{corollary}

\begin{proof}
At \(z=1\), each factor in \eqref{eq:baseq-product} is
\(1+\zeta+\cdots+\zeta^{q-1}=0\).  Thus the product has a zero of order at least \(m\), and
the derivative argument from \cref{thm:prouhet} applies.
\end{proof}

For \(q=2\) and \(\zeta=-1\), \(u_n=\eps_n\), the factor is \(1-z^{2^j}\), and all the binary
formulas above are recovered.

\section{Worked examples and internal consistency checks}
\label{sec:examples}

\subsection{A table through \(n=15\)}

The columns below cross-check four ostensibly different formulas: digit parity, the valuation
of the central binomial coefficient, and the number of odd entries in a Pascal row.

\begin{center}
\small
\begin{tabular}{@{}rcrrrrr@{}}
\toprule
\(n\) & binary & \(\swt(n)\) & \(\tau(n)\) & \(\eps_n\) &
\(\vtwo\binom{2n}{n}\) & \(O(n)\)\\
\midrule
0  & 0000 & 0 & 0 &  1 & 0 & 1\\
1  & 0001 & 1 & 1 & -1 & 1 & 2\\
2  & 0010 & 1 & 1 & -1 & 1 & 2\\
3  & 0011 & 2 & 0 &  1 & 2 & 4\\
4  & 0100 & 1 & 1 & -1 & 1 & 2\\
5  & 0101 & 2 & 0 &  1 & 2 & 4\\
6  & 0110 & 2 & 0 &  1 & 2 & 4\\
7  & 0111 & 3 & 1 & -1 & 3 & 8\\
8  & 1000 & 1 & 1 & -1 & 1 & 2\\
9  & 1001 & 2 & 0 &  1 & 2 & 4\\
10 & 1010 & 2 & 0 &  1 & 2 & 4\\
11 & 1011 & 3 & 1 & -1 & 3 & 8\\
12 & 1100 & 2 & 0 &  1 & 2 & 4\\
13 & 1101 & 3 & 1 & -1 & 3 & 8\\
14 & 1110 & 3 & 1 & -1 & 3 & 8\\
15 & 1111 & 4 & 0 &  1 & 4 & 16\\
\bottomrule
\end{tabular}
\end{center}

\subsection{Three evaluations at \(n=13\)}

For \(n=13\), the binary digits are at positions \(0,2,3\), so
\[
 \eps_{13}=(1-2)^3=-1,
 \qquad
 \tau(13)=1.
\]
The floor sum gives
\[
 13-\left(\left\lfloor\frac{13}{2}\right\rfloor
 +\left\lfloor\frac{13}{4}\right\rfloor
 +\left\lfloor\frac{13}{8}\right\rfloor\right)
 =13-(6+3+1)=3.
\]
The central binomial formula gives
\(\vtwo\binom{26}{13}=3\).  Finally, row \(13\) of Pascal's triangle has
\(2^3=8\) odd entries.  Each route recovers the same odd parity.

\subsection{A compact verification script}

The following pseudocode verifies the main finite identities up to a chosen bound.  It is not
part of any proof, but it is useful when translating the formulas into a proof assistant or a
computer algebra system.

\begin{lstlisting}[style=pseudo]
for n from 0 to N:
    s := popcount(n)
    eps := (-1)^s
    assert v2(binomial(2*n,n)) == s
    oddPascal := count k in [0,n] with binomial(n,k) odd
    assert oddPascal == 2^s

for m from 0 to M:
    P := sum_{n=0}^{2^m-1} eps(n)*z^n
    assert P == product_{j=0}^{m-1} (1-z^(2^j))
    for r from 0 to m-1:
        assert sum_{n=0}^{2^m-1} eps(n)*n^r == 0
\end{lstlisting}

\section{Formula selection guide and a Lean formalization roadmap}
\label{sec:roadmap}

\subsection{Which formula is best for which purpose?}

\renewcommand{\arraystretch}{1.18}
\begin{longtable}{@{}p{0.22\textwidth}p{0.34\textwidth}p{0.36\textwidth}@{}}
\toprule
Goal & Recommended form & Reason \\
\midrule
\endfirsthead
\toprule
Goal & Recommended form & Reason \\
\midrule
\endhead
Fast evaluation & even--odd recurrence or XOR fold & Linear in the number of bits, constant memory. \\
Arithmetic concealment of bits & \(\vtwo\binom{2n}{n}\) & A single classical integer captures the full weight. \\
Pascal-triangle interpretation & \(O(n)=2^{\swt(n)}\) & Support size in characteristic two. \\
Higher-dimensional Pascal analogue & \(O_r(n)=r^{\swt(n)}\) & Works for every number of variables \(r\). \\
Analytic estimates & \(E(z)=\prod_j(1-z^{2^j})\) & Products expose zeros, scales, and Fourier magnitudes. \\
Exact cancellation & finite-difference operator product & Vanishing moments become immediate. \\
Arbitrary Pascal row & submask product \(R_n(z)\) & Restricts Prouhet cancellation to odd row positions. \\
Finite harmonic analysis & Walsh character & The sequence is one pure Walsh frequency. \\
Spectral correlations & Riesz product and \(\eta\)-recursion & Converts substitution into a measure recursion. \\
Characteristic-two algebra & Artin--Schreier equation & Gives a degree-two algebraic certificate. \\
\bottomrule
\end{longtable}
\renewcommand{\arraystretch}{1}

\subsection{Separation between proved source and proposed extensions}

The repository exposition names the primitive-recursive evaluator \(\lean{tmBitPR}\), its
agreement theorem \(\lean{tmBitPR_eq_thueMorseBit}\), and the exact Pascal-row logarithm
result in section 11.8.  A conservative formalization plan for the additional results would
separate independent themes:

\begin{enumerate}
\item \path{ThueMorseValuation.lean}: Legendre, central-binomial, Catalan, successor, and carry
      formulas.
\item \path{ThueMorseMultinomial.lean}: the Frobenius factorization
      \eqref{eq:multinomial-frob-factor}, the count \(O_r(n)=r^{\swt(n)}\), and the signed
      \(r\)-nomial logarithm.
\item \path{ThueMorseGeneratingProduct.lean}: \(P_m\), coefficient extraction, and the
      polynomial Mahler recurrences.  The analytic infinite product can be layered on after
      the finite algebra is complete.
\item \path{ThueMorseProuhet.lean}: exact vanishing order, polynomial cancellation, operator
      factorization, and first nonzero moments.
\item \path{ThueMorsePascalSupport.lean}: Lucas submask support and the sparse Prouhet theorem.
\item \path{ThueMorseFormalSeriesF2.lean}: the functional equation and Artin--Schreier
      polynomial for \(\Theta\).
\end{enumerate}

The finite algebraic statements are natural first targets: they avoid convergence, complex
analysis, and measure theory, while providing the lemmas needed by the analytic layers.

\subsection{Proof-dependency graph}

The report's core can be organized as the following directed chain:
\[
 \text{binary digits}
 \Longrightarrow
 \begin{cases}
  \text{even--odd recursion},\\
  \text{Lucas/Frobenius support},\\
  \text{finite product }P_m,
 \end{cases}
\]
\[
 P_m
 \Longrightarrow
 \begin{cases}
  \text{Mahler equation},\\
  \text{Fourier product},\\
  \text{order-}\!m\text{ zero at }1,
 \end{cases}
 \Longrightarrow
 \begin{cases}
  \text{Prouhet cancellation},\\
  \text{all finite moments},\\
  \text{finite Riesz products}.
 \end{cases}
\]
Independently, Legendre's formula gives the factorial and central-binomial valuation branch,
while Lucas's theorem gives the Pascal and multinomial support branch.  The sparse Prouhet
theorem is where the product and Pascal branches meet.

\section{Concluding synthesis}
\label{sec:conclusion}

The binomial--logarithm formula from section 11.8 is not an isolated curiosity.  It sits at the
intersection of three structures:

\begin{enumerate}
\item binary digit parity, encoded by \((-1)^{\swt(n)}\);
\item Frobenius factorization in characteristic two, which turns binary digits into polynomial
      support; and
\item exact logarithms of support counts, which recover the digit weight before taking its
      parity.
\end{enumerate}

Once those structures are separated, several families become inevitable.  Replacing a
binomial expansion by an \(r\)-nomial expansion changes \(2^{\swt(n)}\) into
\(r^{\swt(n)}\) without changing the outer parity.  Replacing the full initial interval
\([0,2^m)\) by the odd support of a general Pascal row changes the uniform dyadic finite
difference grid into a sparse grid with steps at the one-bit positions of \(n\).  Replacing
coefficient support by evaluation on the unit circle produces Fourier products and Riesz
measures.  Replacing ordinary coefficients by characteristic-two formal series turns the
finite automaton into an Artin--Schreier equation.

The most reusable principle is therefore not one closed formula but a translation dictionary:
\[
 \boxed{
 \text{binary digits}
 \longleftrightarrow
 \text{Frobenius factors}
 \longleftrightarrow
 \text{Boolean-lattice signs}
 \longleftrightarrow
 \text{dyadic finite differences}.}
\]
The Thue--Morse sequence is the common coefficient system of all four descriptions.  This is
why a parity bit can appear, with no visible binary expansion, in a factorial valuation, a row
of Pascal's triangle, a lacunary product, a Walsh character, a Prouhet partition, or a
singular spectral measure.

\begin{thebibliography}{99}

\bibitem{alloucheshallit1999}
J.-P. Allouche and J. Shallit,
\emph{The ubiquitous Prouhet--Thue--Morse sequence},
in C. Ding, T. Helleseth, and H. Niederreiter (eds.),
\emph{Sequences and Their Applications (SETA '98)},
Springer Series in Discrete Mathematics and Theoretical Computer Science,
Springer, London, 1999, pp. 1--16.

\bibitem{alloucheshallit2003}
J.-P. Allouche and J. Shallit,
\emph{Automatic Sequences: Theory, Applications, Generalizations},
Cambridge University Press, Cambridge, 2003.

\bibitem{alloucheriasatshallit2019}
J.-P. Allouche, S. Riasat, and J. Shallit,
\emph{More infinite products: Thue--Morse and the gamma function},
The Ramanujan Journal \textbf{49} (2019), 115--128;
\url{https://arxiv.org/abs/1709.03398}.

\bibitem{baakegrimm2009}
M. Baake and U. Grimm,
\emph{The singular continuous diffraction measure of the Thue--Morse chain},
Journal of Physics A: Mathematical and Theoretical \textbf{41} (2008), 422001;
\url{https://arxiv.org/abs/0809.0580}.

\bibitem{christol1980}
G. Christol, T. Kamae, M. Mend\`es France, and G. Rauzy,
\emph{Suites alg\'ebriques, automates et substitutions},
Bulletin de la Soci\'et\'e Math\'ematique de France \textbf{108} (1980), 401--419.

\bibitem{fine1947}
N. J. Fine,
\emph{Binomial coefficients modulo a prime},
The American Mathematical Monthly \textbf{54} (1947), 589--592.

\bibitem{kummer1852}
E. E. Kummer,
\emph{\"Uber die Erg\"anzungss\"atze zu den allgemeinen Reciprocit\"atsgesetzen},
Journal f\"ur die reine und angewandte Mathematik \textbf{44} (1852), 93--146.

\bibitem{mahler1929}
K. Mahler,
\emph{Arithmetische Eigenschaften der L\"osungen einer Klasse von Funktionalgleichungen},
Mathematische Annalen \textbf{101} (1929), 342--366.

\bibitem{prouhet1851}
E. Prouhet,
\emph{M\'emoire sur quelques relations entre les puissances des nombres},
Comptes rendus de l'Acad\'emie des sciences de Paris \textbf{33} (1851), 225.

\bibitem{queffelec2010}
M. Queff\'elec,
\emph{Substitution Dynamical Systems -- Spectral Analysis},
2nd ed., Lecture Notes in Mathematics 1294, Springer, Berlin, 2010.

\bibitem{reshetnikovrepo}
V. Reshetnikov,
\emph{ProveIt: Analysis/FabiusFunction}, active Lean formalization repository,
\url{https://github.com/VladimirReshetnikov/ProveIt/tree/main/Analysis/FabiusFunction},
repository snapshot 25 August 2026.

\bibitem{reshetnikovexposition}
V. Reshetnikov,
\emph{The Fabius Function and Rvachev's Up-Function: Exact arithmetic, probability, Fourier
analysis, and corrected sharp asymptotics},
LaTeX source in the \texttt{ProveIt} repository, especially Section 11.8,
\url{https://github.com/VladimirReshetnikov/ProveIt/blob/main/Analysis/FabiusFunction/docs/Fabius_Function_and_Rvachev_Up/Fabius_Function_and_Rvachev_Up.tex}.

\bibitem{wright1959}
E. M. Wright,
\emph{Prouhet's 1851 solution of the Tarry--Escott problem of 1910},
The American Mathematical Monthly \textbf{66} (1959), 199--201.

\end{thebibliography}

\end{document}
